\chapter{TMJ Surgery}

\section*{Introduction \& Clinical Indications}
TMJ surgery refers to a variety of surgical interventions aimed at addressing temporomandibular joint disorders (TMDs) that have not responded to conservative treatments. The temporomandibular joint (TMJ) connects the mandible to the temporal bone of the skull, facilitating essential functions such as chewing, speaking, and swallowing. Surgical interventions are indicated when patients experience severe pain, dysfunction, or structural abnormalities that significantly impair their quality of life. The procedures range from minimally invasive techniques, such as arthrocentesis and arthroscopy, to more extensive open surgeries like arthrotomy, condylectomy, and total joint replacement. The primary objectives of TMJ surgery include pain relief, restoration of normal jaw function, and correction of anatomical derangements.

\begin{itemize}
  \item Temporomandibular joint arthroplasty
  \item TMJ arthroscopy with lysis and lavage
  \item Arthrocentesis of the temporomandibular joint
  \item Open TMJ surgery (e.g., disc repositioning, condylectomy)
  \item Total temporomandibular joint replacement
  \item Mandibular condylar surgery
\end{itemize}

The surgical principles of TMJ surgery focus on addressing the underlying anatomical or pathological issues contributing to TMDs. Common conditions include internal derangement, degenerative joint disease, and trauma. The surgical approach aims to restore normal joint mechanics, alleviate pain, and improve function. 

For instance, arthrocentesis involves the hydraulic lavage of the joint to remove inflammatory mediators and debris, thereby reducing intracapsular pressure. Arthroscopy allows for direct visualization and treatment of internal derangements, while open surgery provides access for more extensive reconstruction, such as disc repair or total joint replacement. The overarching goal is to enhance the patient's quality of life by correcting mechanical dysfunction and alleviating pain.

The history of TMJ surgery dates back to the late 19th century, with early attempts focusing on open procedures for ankylosis and trauma. A significant advancement occurred in 1975 when Dr. Hitoshi Ohnishi performed the first TMJ arthroscopy, revolutionizing the treatment of internal derangements. This technique was further refined by Dr. Leon D. McCain in the 1980s, who introduced advanced arthroscopic techniques for disc repositioning.

In the 1990s, arthrocentesis gained popularity as a less invasive option for treating acute closed lock and inflammatory conditions. Concurrently, total TMJ replacement evolved, transitioning from custom-fabricated prostheses to standardized, biocompatible systems in the late 20th century. These advancements have transformed TMJ surgery from a last resort to a spectrum of effective treatments.

TMJ surgery is indicated for various clinical scenarios, including:

\begin{itemize}
  \item Chronic, intractable TMJ pain impacting daily activities and quality of life.
  \item Persistent jaw locking, either open or closed, unresponsive to conservative management.
  \item Severe trismus due to internal derangement or ankylosis.
  \item Diagnosed internal derangement, such as non-reducing disc displacement with severe pain.
  \item Advanced degenerative joint disease causing significant pain and functional impairment.
  \item TMJ ankylosis, severely restricting jaw movement.
  \item Tumors or cysts affecting the TMJ requiring surgical excision.
  \item Trauma to the joint structures necessitating surgical reconstruction.
  \item Failed previous TMJ surgeries requiring revision.
\end{itemize}

TMJ surgery is primarily considered a secondary or tertiary intervention. Most patients undergo a comprehensive course of conservative therapies, including physical therapy, pharmacotherapy, and occlusal splint therapy, before surgical options are considered. Surgery is reserved for those with severe, persistent symptoms unresponsive to conservative management.

In rare cases, such as severe ankylosis or aggressive tumors, surgery may be the primary treatment option. The decision to proceed with surgery depends on the severity of the condition and the failure of conservative measures.

\section*{Relevant Surgical Anatomy}
The temporomandibular joint is a complex synovial joint composed of the mandibular condyle, the articular disc, and the glenoid fossa of the temporal bone. The joint is surrounded by a fibrous capsule and is innervated by the auriculotemporal nerve, a branch of the mandibular division of the trigeminal nerve (CN V3). The arterial supply is primarily from the superficial temporal artery and the maxillary artery, while venous drainage occurs through the pterygoid plexus. Lymphatic drainage is via the preauricular and submandibular lymph nodes.

Key anatomical landmarks include:

\begin{itemize}
  \item McBurney's point: Not applicable in TMJ surgery but important for other abdominal surgeries.
  \item Calot's triangle: Not directly relevant but highlights the importance of understanding anatomical relationships in surgery.
\end{itemize}

Understanding these anatomical relationships is crucial for avoiding complications during surgical procedures.

\section*{Preoperative Workup \& Preparation}
Preoperative preparation is essential for optimizing surgical outcomes. Key components include:

\begin{itemize}
  \item \textbf{Diagnostic Imaging:} MRI is used to evaluate soft tissue structures, while CT scans assess bone morphology and degenerative changes.
  
  \item \textbf{Dental Evaluation:} A thorough dental assessment is performed to evaluate occlusion and dental health.

  \item \textbf{Medical Clearance:} Patients undergo a comprehensive medical evaluation, including blood tests and ECG, to ensure fitness for surgery.

  \item \textbf{Medication Review:} Patients are advised to discontinue blood-thinning medications prior to surgery.

  \item \textbf{NPO Status:} Patients must fast for 6-8 hours before surgery to prevent aspiration.

  \item \textbf{Prophylactic Antibiotics:} Administered shortly before the procedure to reduce infection risk.

  \item \textbf{Informed Consent:} A detailed discussion of the procedure, risks, and benefits is conducted to ensure informed consent.
\end{itemize}

Careful consideration of contraindications is crucial for patient safety.

\textbf{Absolute Contraindications:}
\begin{itemize}
  \item Active infection in the TMJ or surrounding tissues.
  \item Uncontrolled systemic medical conditions.
  \item Severe coagulopathy that cannot be corrected.
  \item Pregnancy, particularly in the first trimester.
  \item Malignancy in the TMJ region requiring non-surgical treatment.
\end{itemize}

\textbf{Relative Contraindications and Precautions:}
\begin{itemize}
  \item Unrealistic patient expectations regarding outcomes.
  \item Significant psychiatric disorders complicating recovery.
  \item Active substance abuse.
  \item Previous radiation therapy to the head and neck.
  \item Severe osteoporosis affecting implant stability.
  \item History of keloid formation.
  \item Non-compliance with postoperative instructions.
\end{itemize}

\section*{Surgical Technique \& Procedure}
The surgical technique for TMJ surgery varies based on the procedure performed. All surgeries are conducted under sterile conditions, typically in an operating room, with general anesthesia being the most common choice.

\begin{itemize}
  \item \textbf{Anesthesia and Positioning:} General anesthesia is preferred for arthroscopy and open surgery. The patient is positioned supine, with the head turned to the contralateral side. The surgical area is prepped and draped.

  \item \textbf{Arthrocentesis:}
    \begin{enumerate}
      \item Identify anatomical landmarks to locate the superior joint space.
      \item Insert two needles into the superior joint compartment.
      \item Infuse sterile saline under pressure to distend the joint and flush out debris.
      \item Manipulate the jaw to lyse adhesions and free a displaced disc.
      \item Allow fluid to drain and remove the needles.
    \end{enumerate}

  \item \textbf{TMJ Arthroscopy:}
    \begin{enumerate}
      \item Make a small incision in the preauricular region and insert a trocar and cannula.
      \item Introduce an arthroscope for visualization of joint structures.
      \item Create a second working portal for miniature instruments.
      \item Inspect the joint for disc position and perform therapeutic procedures as needed.
      \item Close the incisions with fine sutures.
    \end{enumerate}

  \item \textbf{Open TMJ Surgery:}
    \begin{enumerate}
      \item Make a preauricular incision for adequate exposure.
      \item Dissect through tissues, identifying and protecting the facial nerve.
      \item Expose the joint capsule and perform necessary procedures, such as disc repositioning or total joint replacement.
      \item Repair the joint capsule and close the incision in layers.
      \item Place a small drain if necessary.
    \end{enumerate}
\end{itemize}

\section*{Complications \& Risks}
TMJ surgery carries inherent risks, categorized into general surgical risks and procedure-specific complications.

\textbf{General Surgical Risks:}
\begin{itemize}
  \item Anesthesia complications, including allergic reactions and respiratory issues.
  \item Bleeding, which may require transfusion.
  \item Infection, potentially necessitating further treatment.
  \item Scarring, which can be hypertrophic or keloidal.
\end{itemize}

\textbf{Procedure-Specific Risks:}
\begin{itemize}
  \item Facial nerve injury, leading to temporary or permanent weakness.
  \item Ear canal injury, potentially affecting hearing.
  \item Salivary gland injury, resulting in fistula formation.
  \item Trismus or joint stiffness post-surgery.
  \item Malocclusion due to changes in joint anatomy.
  \item Persistent pain or worsening symptoms.
  \item Recurrence of pathology requiring further intervention.
  \item Hardware failure in total joint replacement.
  \item Heterotopic ossification around the joint.
  \item Numbness from sensory nerve injury.
\end{itemize}

\section*{Postoperative Care \& Recovery}
Postoperative recovery varies based on the procedure performed. 

In the immediate postoperative period, patients are monitored in the Post-Anesthesia Care Unit (PACU) for vital signs and pain management. Pain control is initiated with analgesics, and patients may be discharged the same day or require an overnight stay.

During the first 24-48 hours, swelling and discomfort are common. Ice packs can help manage these symptoms. A soft diet is recommended to minimize jaw movement, and oral medications are prescribed. Gentle jaw exercises may begin early to prevent stiffness.

In the first 1-2 weeks, the focus remains on pain control and gradual increase in jaw mobility. Patients are advised to avoid hard or chewy foods and strenuous activities. Long-term recovery can extend for several months, with ongoing physical therapy to restore function.

During TMJ surgery, surgeons document various findings, including the condition of the articular disc, joint surfaces, and surrounding tissues. Pathology reports on resected tissues provide definitive diagnoses, confirming the presence of degenerative changes, inflammation, or tumors. 

In cases of total joint replacement, imaging studies postoperatively assess implant positioning and stability, ensuring proper healing.

A successful postoperative course is characterized by significant pain reduction, improved jaw function, and a return to normal activities. Patients should experience a noticeable improvement in quality of life, with fewer episodes of locking or clicking. 

The timeline for recovery varies, but most patients achieve functional improvement within weeks to months, depending on the procedure performed and adherence to postoperative care.

Postoperative care is critical for ensuring optimal recovery. 

\begin{itemize}
  \item \textbf{Regular Clinical Visits:} Scheduled appointments to monitor healing and assess jaw function.
  
  \item \textbf{Suture Removal:} Non-resorbable sutures are typically removed 7-10 days post-surgery.

  \item \textbf{Physical Therapy (PT):} Ongoing therapy to improve range of motion and strength.

  \item \textbf{Dental Follow-up:} Regular check-ups to monitor occlusion and oral health.

  \item \textbf{Imaging Studies:} Postoperative imaging may be performed to assess for complications.

  \item \textbf{Medication Adjustments:} Pain medications may be adjusted as recovery progresses.

  \item \textbf{Activity Resumption:} Gradual return to normal activities guided by the surgeon.

  \item \textbf{Warning Signs:} Patients should be educated on signs that require immediate medical attention, such as increasing pain, swelling, or fever.
\end{itemize}

\section*{Outcomes \& Prognosis}
Temporomandibular disorders (TMDs) affect an estimated 5-12\% of the adult population, with a higher prevalence in women, particularly those of childbearing age. However, only a small subset, approximately 5-10\%, will require surgical intervention after failing conservative management. 

Arthrocentesis and arthroscopy are the most commonly performed surgical procedures, reflecting their effectiveness for many internal derangements. Open joint surgery and total joint replacement are reserved for more severe cases, with lower incidence rates. Mortality directly attributable to TMJ surgery is rare, but morbidity can vary based on the procedure's invasiveness.

The prognosis for TMJ surgery is generally favorable, with success rates varying based on the procedure and underlying pathology. 

For arthrocentesis, success rates for pain relief and improved function range from 70-90\%. TMJ arthroscopy shows similar outcomes, with success rates between 70-85\%. Open joint procedures can achieve comparable results but carry higher risks.

Total joint replacement typically yields excellent functional outcomes, with significant pain reduction in 85-95\% of patients. Recurrence of symptoms can occur, particularly with less invasive procedures. Prognostic factors influencing outcomes include the duration of symptoms, psychological comorbidities, and patient compliance with postoperative care.

\section*{References}
\begin{enumerate}
  \item MedlinePlus. TMJ Surgery. U.S. National Library of Medicine; 2025.
  \item Hupp JR, Ellis E, Tucker MR. Contemporary Oral and Maxillofacial Surgery. 7th ed. Elsevier; 2019.
  \item American Association of Oral and Maxillofacial Surgeons. Parameters of Care: Clinical Practice Guidelines for Oral and Maxillofacial Surgery. 6th ed. AAOMS; 2017.
  \item Nitzan DW, Dolwick MF, Marmary Y, et al. The effect of arthrocentesis on patient's symptoms and joint pathology. J Oral Maxillofac Surg. 1991;49(11):1163-1169.
\end{enumerate}

\clearpage
